\section*{Chapter \ref{ch:g7:plate-tectonics} --- Plate Tectonics: The Moving Puzzle}

\begin{solution}{exo:g7:plate-tectonics:1}
\omterm{def:g7:plate-tectonics:continental-drift}{Continental drift} is the idea that the \omterm{def:g3:continents-and-oceans:continent}{continents} move slowly across the
Earth's surface, so that their positions have changed over the ages.
Wegener's clues: the fit of the coastlines (South America and Africa); the
same fossils on \omterm{def:g3:continents-and-oceans:continent}{continents} now separated by \omterm{def:g2:water-on-earth:ocean}{oceans}; the same \omterm{def:g3:stones-pebbles-sand:rock}{rocks} and old
\omterm{def:g1:landscapes-around-us:landform}{mountain} chains on either side of an \omterm{def:g2:water-on-earth:ocean}{ocean}; traces of ancient \omterm{def:g5:weather-and-climate:climate}{climates} out
of place (coal in Antarctica, glacial deposits in hot, dry Africa).
\end{solution}

\begin{solution}{exo:g7:plate-tectonics:2}
Divergent: the \omterm{def:g7:plate-tectonics:plate}{plates} move apart. Convergent: they move towards each other.
Transform: they slide past each other, side by side.
\end{solution}

\begin{solution}{exo:g7:plate-tectonics:3}
For example, all around the Pacific \omterm{def:g2:water-on-earth:ocean}{Ocean} (Japan, Indonesia, the Andes),
and the belt running from the Mediterranean to the Himalaya. Almost no
large \omterm{def:g4:volcanoes-earthquakes-first-look:earthquake}{earthquakes} in the middle of Africa or of Australia (also of North
America east of the Rocky \omterm{def:g1:landscapes-around-us:landform}{Mountains}, or of Brazil).
\end{solution}

\begin{solution}{exo:g7:plate-tectonics:4}
An animal of \omterm{def:g2:water-on-earth:salt-fresh}{fresh water} could not cross thousands of kilometres of salt
\omterm{def:g2:water-on-earth:ocean}{ocean}. Finding it on both sides means that the two places were once joined
by land, in one \omterm{def:g3:continents-and-oceans:continent}{continent}, when \emph{Mesosaurus} lived.
\end{solution}

\begin{solution}{exo:g7:plate-tectonics:5}
Coal forms from the plants of warm, wet, swampy forests: the \omterm{def:g5:weather-and-climate:climate}{climate} was
warm and humid. Such forests cannot grow near the \omterm{def:g3:continents-and-oceans:poles}{South Pole} today, so
Antarctica must once have been much closer to the \omterm{def:g3:continents-and-oceans:equator}{equator}: it has moved.
\end{solution}

\begin{solution}{exo:g7:plate-tectonics:6}
Mid-Atlantic Ridge: divergent. Andes: convergent (an \omterm{def:g2:water-on-earth:ocean}{ocean} \omterm{def:g7:plate-tectonics:plate}{plate} sinks
under a \omterm{def:g3:continents-and-oceans:continent}{continent}). San Andreas fault: transform. Himalaya: convergent (two
\omterm{def:g3:continents-and-oceans:continent}{continents} collide). East African Rift: divergent (a \omterm{def:g3:continents-and-oceans:continent}{continent} splitting).
\end{solution}

\begin{solution}{exo:g7:plate-tectonics:7}
$2.5 \times 10 = 25$ kilometres per million years; in 60 million years:
$25 \times 60 = \qty{1500}{km}$.
\end{solution}

\begin{solution}{exo:g7:plate-tectonics:8}
$\qty{7}{m} = \qty{700}{cm}$, and $700 \div 9 \approx 78$: about
\qty{78}{cm} per year, some 30 times the average of \qty{2.5}{cm} per year.
The \omterm{def:g7:plate-tectonics:plate}{plates} do not slide apart smoothly at the surface: the strain builds up
for years and is released in short bursts, when the cracks open by metres
at once. Over long periods, quiet years and bursts average out to a few
centimetres per year.
\end{solution}

\begin{solution}{exo:g7:plate-tectonics:9}
He could not say what force moves the \omterm{def:g3:continents-and-oceans:continent}{continents}, and the motion he
imagined (\omterm{def:g3:continents-and-oceans:continent}{continents} ploughing through the \omterm{def:g2:water-on-earth:ocean}{ocean} floor) was impossible.
In \omterm{def:g7:plate-tectonics:plate-tectonics}{plate tectonics}, the \omterm{def:g3:continents-and-oceans:continent}{continents} do not plough through anything: they are
carried, with the \omterm{def:g2:water-on-earth:ocean}{ocean} floor around them, by \omterm{def:g7:plate-tectonics:plate}{plates} that move as whole
pieces of the Earth's outer shell.
\end{solution}

\begin{solution}{exo:g7:plate-tectonics:10}
Highest at the ridge, deepest at the trench. New \omterm{def:g7:plate-tectonics:plate}{plate} is made at the
ridge, where hot \omterm{def:g3:stones-pebbles-sand:rock}{rock} rises from the \omterm{def:g6:inside-the-earth:layers}{mantle}; the \omterm{def:g2:water-on-earth:ocean}{ocean} \omterm{def:g7:plate-tectonics:plate}{plate} goes back into
the \omterm{def:g6:inside-the-earth:layers}{mantle} at the trench, where it sinks beneath the \omterm{def:g3:continents-and-oceans:continent}{continent}.
\end{solution}

\begin{solution}{exo:g7:plate-tectonics:11}
$14.2 \div 2.2 \approx 6.5$: about 6 times faster (6.45, rounded). In
kilometres per million years the speeds are 142 and 22. Times for
\qty{1000}{km}: $1000 \div 142 \approx 7$ million years for the East
Pacific Rise; $1000 \div 22 \approx 45$ million years for the North
Atlantic ridge.
\end{solution}

\begin{solution}{exo:g7:plate-tectonics:12}
(a) $6472 - 4578 = 1894$ \omterm{def:g4:volcanoes-earthquakes-first-look:earthquake}{earthquakes}. (b) $4578 \div 7031 \approx 0.651$,
that is 65\,\%. (c) The line on the \omterm{def:g2:drawing-to-map:map}{map} marks the trench, where the \omterm{def:g2:water-on-earth:ocean}{ocean}
\omterm{def:g7:plate-tectonics:plate}{plate} starts to sink. The \omterm{def:g7:plate-tectonics:plate}{plate} does not go straight down: it sinks
slanting beneath the other \omterm{def:g7:plate-tectonics:plate}{plate}, deeper and deeper as it goes under it.
\omterm{def:g4:volcanoes-earthquakes-first-look:earthquake}{Earthquakes} happen all along the sinking \omterm{def:g7:plate-tectonics:plate}{plate}, so many of them lie under
the overriding \omterm{def:g7:plate-tectonics:plate}{plate}, on the same side of the trench, tens or hundreds of
kilometres from the line. A \omterm{def:g7:plate-tectonics:plate-boundary}{plate boundary} is a zone, not a thin line.
\end{solution}

\begin{solution}{pb:g7:plate-tectonics:1}
\textbf{1.} \emph{Mesosaurus}, \emph{Cynognathus} and \emph{Glossopteris}
(the band of \emph{Lystrosaurus} joins Africa, India and Antarctica).

\textbf{2.} It lived in \omterm{def:g2:water-on-earth:salt-fresh}{fresh water}: it could not cross a salt \omterm{def:g2:water-on-earth:ocean}{ocean}
thousands of kilometres wide.

\textbf{3.} The bulge of north-east Brazil fits into the hollow of the
Gulf of Guinea, on Africa's west \omterm{def:g1:landscapes-around-us:water}{coast}.

\textbf{4.} A \omterm{def:g7:plate-tectonics:plate-boundary}{divergent boundary}, the Mid-Atlantic Ridge: the South
American \omterm{def:g7:plate-tectonics:plate}{plate} moves west and the African \omterm{def:g7:plate-tectonics:plate}{plate} east, away from each other.

\textbf{5.} The \omterm{def:g1:landscapes-around-us:water}{coasts} are on opposite sides of the \omterm{def:g6:latitude-longitude:longitude}{prime meridian}:
$43.2 + 14.4 = \ang{57.6}$.

\textbf{6.} $57.6 \times 102 = 5875.2$: about \qty{5880}{km}.

\textbf{7.} The \omterm{def:g6:latitude-longitude:longitude}{meridians} come closer together as one goes from the
\omterm{def:g3:continents-and-oceans:equator}{equator} towards a \omterm{def:g3:continents-and-oceans:poles}{pole}. The \omterm{def:g6:latitude-longitude:latitude}{parallel} at \ang{23} South is a smaller circle
than the \omterm{def:g3:continents-and-oceans:equator}{equator}, so each of its 360 degrees is shorter.

\textbf{8.} About half the width: $5880 \div 2 = \qty{2940}{km}$.

\textbf{9.} $3.4 \times 10 = 34$ kilometres per million years.

\textbf{10.} $5880 \div 34 \approx 173$: about 170 million years.

\textbf{11.} Yes: 170 million years lies between 100 and 200 million years,
and it comes after the start of the break-up of \omterm{def:g7:plate-tectonics:supercontinent}{Pangaea}, 225 to 200 million
years ago. A simple calculation on a \omterm{def:g2:drawing-to-map:map}{map} gives the right order of
magnitude.

\textbf{12.} For example: the speed has not been the same all the time;
the \omterm{def:g3:continents-and-oceans:continent}{continents} stretched and thinned before they split, so today's \omterm{def:g1:landscapes-around-us:water}{coasts}
are not exactly the line along which they broke apart.
\end{solution}
